% =============================================================================
% Хавсралт Б: Eval set-ийн бичлэгийн бүтэц
% =============================================================================

\chapter{Eval set-ийн бичлэгийн бүтэц}
\label{AppendixB}
\pagecolor{white}

Eval set нь \tsql{evalset/dev.jsonl} (300) ба \tsql{evalset/test.jsonl} (200) файлд, мөр бүр нэг JSON бичлэг
хэлбэрээр хадгалагдана. Хувилбар бүрийн hash нь \tsql{evalset/MANIFEST.json}-д бичигдэх ба run бүрийн бичлэгт хамт
хадгалагдана. Талбаруудыг Хүснэгт~\ref{tab:eval-fields}-д тайлбарлав.

\begin{table}[H]
\centering\small
\caption{Eval set-ийн бичлэгийн талбарууд}
\label{tab:eval-fields}
\begin{tabularx}{\textwidth}{p{3.6cm}X}
\toprule
Талбар & Утга \\
\midrule
\tsql{id}, \tsql{family} & Асуултын ID ба хэв загварын бүлэг (хуваалт бүлгээр) \\
\tsql{domain} & applications, loans, payments, collection, customers, deposits, gl\_trends, cross \\
\tsql{category} & answerable, unanswerable, ambiguous \\
\tsql{difficulty} & easy, medium, hard (хариулах боломжтой асуултад) \\
\tsql{question\_mn} / \tsql{\_latin} / \tsql{\_en} / \tsql{\_mixed} & Кирилл, латин галиг, англи, холимог хувилбар \\
\tsql{gold\_sql}, \tsql{alt\_sql} & Лавлах SQL ба орлох (тэнцүү) SQL-үүд (\tsql{t2s} схем дээр) \\
\tsql{order\_matters} & Үр дүнгийн эрэмбэ чухал эсэх \\
\tsql{interpretations}, \tsql{clarification\_mn} & Хоёрдмол асуултын тайлбарууд ба тодруулах асуулт \\
\tsql{unanswerable\_reason} & pii, missing\_data, non\_sql, external\_knowledge, future, write\_operation \\
\tsql{tags}, \tsql{contains\_value}, \tsql{values}, \tsql{glossary\_terms} & Задаргааны шошго, асуултад гарсан утга, нэр томьёо \\
\bottomrule
\end{tabularx}
\end{table}

\begin{lstlisting}[basicstyle=\ttfamily\scriptsize, caption={Бичлэгийн жишээ (dev, A001; уншихад хялбар болгож мөр таслав, зарим талбарыг орхив)}, breaklines=true, showstringspaces=false, language={}]
{"id": "A001", "family": "A-f01", "domain": "applications",
 "category": "answerable", "difficulty": "easy",
 "question_mn": "2025 онд нийт хэдэн зээлийн хүсэлт ирсэн бэ?",
 "question_latin": "2025 ond niit heden zeeliin hvselt irsen be?",
 "question_en": "How many loan applications were received in 2025 in total?",
 "gold_sql": "SELECT count(*) FROM t2s.loan_applications
              WHERE submitted_date BETWEEN '2025-01-01' AND '2025-12-31'",
 "order_matters": false, "tags": ["aggregation", "date"]}
\end{lstlisting}

Run бүрийн бичлэг дараах файлд хадгалагдана:
\begin{center}
\tsql{results/runs/<split>/<ID>/<загвар>\_\_<хувилбар>\_\_s<seed>.jsonl}
\end{center}
Бичлэг нь асуултын ID, үүсгэсэн SQL, үр дүнгийн hash, зөв эсэх (EX, хатуу EX, TS-EX), төлөв, алдааны шат (extract,
guard, exec, empty, refusal), итгэл, нэр дэвшигчдийн hash, retrieval-ийн мэдээлэл, agent-ийн алхам, хугацаа, токен,
eval set ба кодын хувилбарыг агуулна.
